\documentclass[10pt]{article}

% Load the custom package
\usepackage{Mathdoc}

\begin{document}


%% Numéro de séquence %% Titre de la séquence %%
\renewcommand{\centerhead}{Fractions -  Devoir Bilan}

%% Spaces Commands %%
\renewcommand{\baselinestretch}{1}
\setlength{\parindent}{0pt}

\entetedevoirs{30}

\begin{center}
  \duree{1 heures} \coefficient{2} \calculatrice{0}
\end{center}
\begin{multicols}{2}
  \begin{exercicedevoir}[5]
    \textbf{Mettre les fractions au même dénominateur \\(Comme dans
    l'exemple 1).}
    \begin{enumerate}
    \item $\dfrac{1}{3} = \dfrac{5 \times 1}{5 \times 3}
      = $$\dfrac{5}{15}$
    \item $\dfrac{5}{9} =
      \dfrac{\phantom{00000000000000}}{\phantom{00000000000000}}
      = $$\dfrac{\phantom{0000}}{72}$
    \item $5 =
      \dfrac{\phantom{00000000000000}}{\phantom{00000000000000}} =
      \dfrac{\phantom{0000}}{7}$
    \item $\dfrac{7}{10} =
      \dfrac{\phantom{00000000000000}}{\phantom{00000000000000}}
      = $$\dfrac{\phantom{0000}}{30}$
    \item $\dfrac{6}{7} =
      \dfrac{\phantom{00000000000000}}{\phantom{00000000000000}}
      = $$\dfrac{\phantom{0000}}{42}$
    \item $\dfrac{5}{8} =
      \dfrac{\phantom{00000000000000}}{\phantom{00000000000000}}
      = $$\dfrac{\phantom{0000}}{80}$
    \item $1 =
      \dfrac{\phantom{00000000000000}}{\phantom{00000000000000}} =
      \dfrac{\phantom{0000}}{7}$
    \item $\dfrac{1}{5} =
      \dfrac{\phantom{00000000000000}}{\phantom{00000000000000}}
      = $$\dfrac{\phantom{0000}}{20}$
    \item $\dfrac{7}{9} =
      \dfrac{\phantom{00000000000000}}{\phantom{00000000000000}}
      = $$\dfrac{\phantom{0000}}{18}$
    \item $\dfrac{1}{8} =
      \dfrac{\phantom{00000000000000}}{\phantom{00000000000000}}
      = $$\dfrac{\phantom{0000}}{48}$
    \end{enumerate}
  \end{exercicedevoir}
\vspace{-2cm}
\begin{exercicedevoir}[5]
  \textbf{Simplifier les fractions suivantes \\(Comme dans l'exemple
  1).}
  \begin{enumerate}
  \item $ \dfrac{40}{64} = \dfrac{\cancel{8} \times
    5}{\cancel{8}\times 8} = \dfrac{5}{8} $
  \item $ \dfrac{5}{30} = \dfrac{\phantom{00000000000000}}{} =
    \dfrac{\phantom{0000}}{} $
  \item $ \dfrac{12}{42} = \dfrac{\phantom{00000000000000}}{} =
    \dfrac{\phantom{0000}}{} $
  \item $ \dfrac{2}{20} = \dfrac{\phantom{00000000000000}}{} =
    \dfrac{\phantom{0000}}{} $
  \item $ \dfrac{42}{48} = \dfrac{\phantom{00000000000000}}{} =
    \dfrac{\phantom{0000}}{} $
  \item $ \dfrac{56}{80} = \dfrac{\phantom{00000000000000}}{} =
    \dfrac{\phantom{0000}}{} $
  \item $ \dfrac{9}{18} = \dfrac{\phantom{00000000000000}}{} =
    \dfrac{\phantom{0000}}{} $
  \item $ \dfrac{20}{24} = \dfrac{\phantom{00000000000000}}{} =
    \dfrac{\phantom{0000}}{} $
  \item $ \dfrac{90}{100} = \dfrac{\phantom{00000000000000}}{} =
    \dfrac{\phantom{0000}}{} $
  \item $ \dfrac{18}{30} = \dfrac{\phantom{00000000000000}}{} =
    \dfrac{\phantom{0000}}{} $
  \end{enumerate}
\end{exercicedevoir}
\end{multicols}

\begin{exercicedevoir}
\textbf{Décomposer les fractions suivantes comme la somme d'un entiers
et d'une fraction décimale. \\ (Comme dans l'exemple 1)}
\begin{multicols}{2}
  \begin{enumerate}
  \item $ \dfrac{31}{10} = 3 + \dfrac{1}{10} $  $\quad$ car $\quad$ $31 = 3 \times 10 +
    1$
  \item $ \dfrac{3}{2} = \ldots\ldots +
    \dfrac{\ldots\ldots}{\ldots\ldots} $ $\quad$ car \dotfill
  \item $ \dfrac{12}{5} = \ldots\ldots +
    \dfrac{\ldots\ldots}{\ldots\ldots} $$\quad$ car \dotfill
  \item $ \dfrac{5}{4} = \ldots\ldots +
    \dfrac{\ldots\ldots}{\ldots\ldots} $$\quad$ car \dotfill
  \item $ \dfrac{9}{4} = \ldots\ldots +
    \dfrac{\ldots\ldots}{\ldots\ldots} $$\quad$ car \dotfill
  \item $ \dfrac{16}{8} = \ldots\ldots +
    \dfrac{\ldots\ldots}{\ldots\ldots} $$\quad$ car \dotfill
  \item $ \dfrac{22}{10} = \ldots\ldots +
    \dfrac{\ldots\ldots}{\ldots\ldots} $$\quad$ car \dotfill
  \item $ \dfrac{18}{5} = \ldots\ldots +
    \dfrac{\ldots\ldots}{\ldots\ldots} $$\quad$ car \dotfill
  \item $ \dfrac{4}{4} = \ldots\ldots +
    \dfrac{\ldots\ldots}{\ldots\ldots} $$\quad$ car \dotfill
  \item $ \dfrac{13}{5} = \ldots\ldots +
    \dfrac{\ldots\ldots}{\ldots\ldots} $$\quad$ car \dotfill
  \end{enumerate}
\end{multicols}
\end{exercicedevoir}

\begin{multicols}{2}
  \begin{exercicedevoir}[5]
    \textbf{Comparer les fractions suivantes.  \\(Comme dans l'exemple
    1)}
    \begin{enumerate}
    \item $\dfrac{37}{77}>\dfrac{3}{7}$ $\quad$ car $\quad$
      $\dfrac{3}{7}=\dfrac{33}{77}$ $\quad$ (en multipliant par 11)
    \item $\dfrac{1}{2}\quad\ldots\quad\dfrac{10}{22}$ $\quad$ car $\quad$ $\dfrac{1}{2}=\dfrac{\ldots\ldots}{22}$
    \item $\dfrac{1}{5}\quad\ldots\quad\dfrac{10}{40}$ $\quad$ car $\quad$ $\dfrac{1}{5}=\dfrac{\ldots\ldots}{40}$
    \item $\dfrac{3}{4}\quad\ldots\quad\dfrac{7}{8}$ $\quad$ car $\quad$ $\dfrac{3}{4}=\dfrac{\ldots\ldots}{8}$
    \item $\dfrac{9}{99}\quad\ldots\quad\dfrac{1}{9}$ $\quad$ car $\quad$ $\dfrac{1}{9}=\dfrac{\ldots\ldots}{99}$
    \item $\dfrac{4}{7}\quad\ldots\quad\dfrac{11}{14}$ $\quad$ car $\quad$ $\dfrac{4}{7}=\dfrac{\ldots\ldots}{14}$
    \item $\dfrac{5}{7}\quad\ldots\quad\dfrac{37}{56}$ $\quad$ car $\quad$ $\dfrac{5}{7}=\dfrac{\ldots\ldots}{56}$
    \item $\dfrac{11}{18}\quad\ldots\quad\dfrac{5}{6}$ $\quad$ car $\quad$ $\dfrac{5}{6}=\dfrac{\ldots\ldots}{18}$
    \item $\dfrac{12}{64}\quad\ldots\quad\dfrac{1}{8}$ $\quad$ car $\quad$ $\dfrac{1}{8}=\dfrac{\ldots\ldots}{64}$
    \item $\dfrac{3}{5}\quad\ldots\quad\dfrac{20}{35}$ $\quad$ car $\quad$ $\dfrac{3}{5}=\dfrac{\ldots\ldots}{35}$
    \end{enumerate}
  \end{exercicedevoir}

  \begin{exercicedevoir}[5]
    \textbf{Calculer les multiplications de fractions suivantes.  \\(Comme dans l'exemple
    1)}
    \begin{enumerate}
    \item $\dfrac{6}{7}\times\dfrac{3}{4} =\dfrac{6\times3}{7\times4}=\dfrac{18}{28} $  
    \item $\dfrac{1}{10}\times\dfrac{1}{9} = \dfrac{\phantom{0000000000000000}}{\phantom{0000000000000000}} = \dfrac{\phantom{00000000}}{\phantom{00000000}}$
	\item $\dfrac{2}{3}\times\dfrac{4}{7}  = \dfrac{\phantom{0000000000000000}}{\phantom{0000000000000000}} = \dfrac{\phantom{00000000}}{\phantom{00000000}}$
	\item $\dfrac{4}{5}\times\dfrac{5}{9} = \dfrac{\phantom{0000000000000000}}{\phantom{0000000000000000}} = \dfrac{\phantom{00000000}}{\phantom{00000000}}$
	\item $\dfrac{2}{9}\times\dfrac{3}{5} = \dfrac{\phantom{0000000000000000}}{\phantom{0000000000000000}} = \dfrac{\phantom{00000000}}{\phantom{00000000}}$
	\item $\dfrac{7}{10}\times\dfrac{4}{9} = \dfrac{\phantom{0000000000000000}}{\phantom{0000000000000000}} = \dfrac{\phantom{00000000}}{\phantom{00000000}}$
	\item $\dfrac{1}{8}\times\dfrac{1}{4} = \dfrac{\phantom{0000000000000000}}{\phantom{0000000000000000}} = \dfrac{\phantom{00000000}}{\phantom{00000000}}$
	\item $\dfrac{3}{10}\times\dfrac{6}{7} = \dfrac{\phantom{0000000000000000}}{\phantom{0000000000000000}} = \dfrac{\phantom{00000000}}{\phantom{00000000}}$
	\item $\dfrac{3}{8}\times\dfrac{1}{10} = \dfrac{\phantom{0000000000000000}}{\phantom{0000000000000000}} = \dfrac{\phantom{00000000}}{\phantom{00000000}}$
	\item $8\times\dfrac{7}{10} = \dfrac{\phantom{0000000000000000}}{\phantom{0000000000000000}} = \dfrac{\phantom{00000000}}{\phantom{00000000}}$
    \end{enumerate}
  \end{exercicedevoir}
\end{multicols}

\begin{exercicedevoir}[5][Problème]
\textbf{Au marché, Emma décide d'acheter des fruits pour préparer une
salade de fruits.} \\ \\
Elle achète d'abord des pommes et des bananes : \\
Elle prend $\dfrac{6}{12}$ de kilogramme de pommes et elle prend aussi $\dfrac{4}{8}$ de kilogramme de bananes.
\begin{enumerate}
\item Simplifie les fractions de kilogrammes de pommes et de bananes
achetés par Emma.  \\ \dtf \\ \dtf \\ \dtf
\end{enumerate} 
Ensuite, elle achète des oranges et des poires : \\
Elle prend $\dfrac{3}{5}$ de kilogramme d’oranges et elle prend
$\dfrac{33}{50}$ de kilogramme de poires.
\begin{enumerate}[resume]
\item  Mets ces deux fractions au même dénominateur. \\ \dtf \\ \dtf \\ \dtf
\item Compare la quantité  de fruits en kilogrammes, a-t-elle acheté plus d'oranges ou de poires ? \\ \dtf \\ \dtf \\ \dtf
\end{enumerate} 
\end{exercicedevoir}

\nonewpage
\end{document}

%%% Local Variables:
%%% mode: LaTeX
%%% TeX-master: t
%%% TeX-master: t
%%% End:

